
\subsubsection{5.1 Detection Performance (RQ1, RQ3)}

\textbf{Table 1: Detection Performance on HaluEval}

\begin{table*}[htbp]
\centering
\resizebox{\dimexpr 2\columnwidth+\columnsep\relax}{!}{%
\begin{tabular}{llllllll}
\toprule
Task & AUROC & DeLong p-value & Cohen's d & AUROC\_max & Gate & SelfCheckGPT-NLI & $\Delta$ \\
\midrule
Dialogue & 0.709 & $\approx$ 0 & 0.714 & 0.77 & PASS & 0.48 & +0.229 \\
QA & 0.644 & 1.$29\times 10^{-282}$ & 0.779 & 0.77 & PASS & 0.53 & +0.114 \\
Summarization & 0.530 & 2.$02\times 10^{-13}$ & 0.220 & 0.52 & FAIL & --- & --- \\
\bottomrule
\end{tabular}
}
\end{table*}

\textit{N = 20,000 per task (balanced: 10,000 correct, 10,000 hallucinated). Gate criterion: AUROC > 0.55 with DeLong p < 0.05. Overall gate: 2/3 PASS. Exact ground-truth values: dialogue AUROC = 0.7094, QA AUROC = 0.6437, summarization AUROC = 0.5300. p\_contradictory (proportion of examples where NLI assigns contradiction): dialogue = 0.54, QA = 0.54, summarization = 0.04.}

On dialogue, AUROC = 0.709 indicates that the frozen NLI classifier ranks hallucinated responses above non-hallucinated ones in approximately 71\% of pairs. The DeLong p-value is zero to machine precision, and Cohen's d = 0.714 constitutes a large effect. On QA, AUROC = 0.644 falls 0.006 below the pre-specified 0.65 target from the original hypothesis; Cohen's d = 0.779 is the largest of the three tasks. On summarization, AUROC = 0.530 is near chance, and the theoretical structural ceiling (AUROC\_max = 0.52) indicates the method is near its architectural upper bound on this task.

Both commission tasks (dialogue, QA) exceed the generation-based baseline (SelfCheckGPT-NLI) by substantial margins. The comparison is bounded by the SelfCheckGPT evaluation context (base LLaMA-3-8B) as noted in Section 4.2.

\begin{figure}[htbp]
\centering
\includegraphics[width=\columnwidth]{figures/roc_curves.png}
\caption{ROC curves for all three HaluEval tasks}
\label{fig:roc_curves}
\end{figure}

\textit{Figure 1: ROC curves for NLI-based hallucination detection on HaluEval dialogue, QA, and summarization tasks. AUROC values: dialogue = 0.709, QA = 0.644, summarization = 0.530.}

\begin{figure}[htbp]
\centering
\includegraphics[width=\columnwidth]{figures/score_distributions.png}
\caption{Net-contradiction score distributions}
\label{fig:score_distributions}
\end{figure}

\textit{Figure 2: Net-contradiction score distributions (P(contradiction) $-$ P(entailment)) for hallucinated (label=1) and non-hallucinated (label=0) responses, by task.}

\begin{figure}[htbp]
\centering
\includegraphics[width=\columnwidth]{figures/gate_metrics_comparison.png}
\caption{Gate metrics comparison}
\label{fig:gate_metrics_comparison}
\end{figure}

\textit{Figure 3: AUROC and Cohen's d by task, with gate threshold reference lines.}

\subsubsection{5.2 Mechanism Analysis: Score Ordering (RQ2)}

\textbf{Table 2: Mechanism Verification Statistics}

\begin{table*}[htbp]
\centering
\resizebox{\dimexpr 2\columnwidth+\columnsep\relax}{!}{%
\begin{tabular}{lllllll}
\toprule
Task & KL Divergence & KL > 0.05 & Wilcoxon p & Wilcoxon Pass & Cohen's d & Near-Uniform \% \\
\midrule
Dialogue & 0.2794 & Yes & $\approx$ 0 & Yes & 0.714 & 0.005\% \\
QA & 0.0353 & No & 1.$52\times 10^{-271}$ & Yes & 0.779 & 0.015\% \\
Summarization & 0.3104 & Yes & 2.$07\times 10^{-13}$ & Yes & 0.220 & 0.005\% \\
\bottomrule
\end{tabular}
}
\end{table*}

\textit{KL divergence computed per class from uniform distribution. Wilcoxon rank-sum test: two-sided. Near-uniform threshold: score entropy within 1\% of maximum.}

Wilcoxon rank-sum tests pass on all three tasks: DeBERTa's contradiction scores are stochastically ordered by hallucination label in every task. This confirms that the mechanism --- score differentiation based on NLI training --- is active across the full dataset.

The pre-specified h-m1 gate required KL > 0.05 on all three tasks. This criterion failed on QA (KL = 0.0353), triggering a SELF\_MODIFY outcome in which Wilcoxon and Cohen's d were used as primary mechanism indicators. The QA KL shortfall appears to reflect score distribution compression in short-context QA: both hallucinated and non-hallucinated classes produce narrow score distributions, preserving ordinal separation (Wilcoxon passes; Cohen's d = 0.779, the highest of all tasks) while reducing global distributional spread (KL). The KL threshold of 0.05 appears too strict for short-context tasks; this is noted as a gate criterion design issue (Limitation L6).

Summarization has the second-highest KL (0.310) and a significant Wilcoxon p-value (2.$07\times 10^{-13})$, yet AUROC = 0.530 and Cohen's d = 0.220. This combination indicates that a weak contradiction signal is present in summarization but is insufficient for practical discrimination. The low p\_contradictory for summarization (0.04, versus 0.54 for dialogue and QA) indicates the NLI model rarely assigns contradiction to summarization pairs, consistent with omission-type hallucinations that do not directly contradict the source.

\begin{figure}[htbp]
\centering
\includegraphics[width=\columnwidth]{figures/score_distributions_violin.png}
\caption{Violin plots of score distributions per task}
\label{fig:score_distributions_violin}
\end{figure}

\textit{Figure 4: Violin plots of net-contradiction score distributions for hallucinated and non-hallucinated responses by task.}

\begin{figure}[htbp]
\centering
\includegraphics[width=\columnwidth]{figures/score_separation_boxplot.png}
\caption{Score separation boxplots}
\label{fig:score_separation_boxplot}
\end{figure}

\textit{Figure 5: Box plots of P(contradiction) by hallucination label for each task, with Wilcoxon rank-sum p-values annotated.}

\begin{figure}[htbp]
\centering
\includegraphics[width=\columnwidth]{figures/kl_divergence_summary.png}
\caption{KL divergence per task}
\label{fig:kl_divergence_summary}
\end{figure}

\textit{Figure 6: KL divergence from uniform per task. Dialogue (0.279) and summarization (0.310) exceed the 0.05 threshold; QA (0.035) does not.}

\begin{figure}[htbp]
\centering
\includegraphics[width=\columnwidth]{figures/near_uniform_proportion.png}
\caption{Near-uniform proportion per task}
\label{fig:near_uniform_proportion}
\end{figure}

\textit{Figure 7: Near-uniform score proportion per task (all tasks near 0\%), indicating DeBERTa produces genuinely graded NLI scores across the full dataset.}

\subsubsection{5.3 Structural Ceiling Analysis (RQ4)}

\begin{figure}[htbp]
\centering
\includegraphics[width=\columnwidth]{figures/structural_ceiling.png}
\caption{Structural ceiling analysis}
\label{fig:structural_ceiling}
\end{figure}

\textit{Figure 8: Structural ceiling analysis showing theoretical AUROC upper bound (AUROC\_max) per task alongside observed AUROC.}

The structural ceiling for summarization is AUROC\_max = 0.52, derived from p\_contradictory $\approx$ 0.04: if only approximately 4\% of hallucinated examples produce a contradiction signal, a detector that perfectly identifies those achieves an expected AUROC of approximately 0.52 over the full balanced dataset, consistent with the observed 0.530. For dialogue and QA, AUROC\_max = 0.77, indicating room above the observed 0.709 and 0.644.

HaluEval-Dialogue and HaluEval-QA hallucinations are constructed to contradict the grounding context (substituting correct facts with plausible-but-incorrect alternatives, inventing entities, replacing correct answers with incorrect ones). These are commission-type hallucinations: the response asserts a claim that the context refutes. HaluEval-Summarization hallucinations are abstractive failures --- responses that omit key information or include content not contradicted by, but not fully supported in, the source document. P(contradiction) is expected to be low for these pairs because the hallucinated summary does not contradict the source; it merely fails to capture its meaning. The observed AUROC = 0.530 near the structural ceiling of 0.52 is consistent with this interpretation.
